\documentclass[a4paper, serif, 10pt]{moderncv}

%% moderncv
% style and color
\moderncvstyle{banking}
\moderncvcolor{black}

% renew moderncv command to adapt for CJK colon
\renewcommand*{\cvitem}[3][.25em]{%
  \ifstrempty{#2}{}{\hintstyle{#2}：}{#3}%
  \par\addvspace{#1}}

\renewcommand*{\cvitemwithcomment}[4][.25em]{%
  \savebox{\cvitemwithcommentmainbox}{\ifstrempty{#2}{}{\hintstyle{#2}：}#3}%
  \setlength{\cvitemwithcommentmainlength}{\widthof{\usebox{\cvitemwithcommentmainbox}}}%
  \setlength{\cvitemwithcommentcommentlength}{\maincolumnwidth-\separatorcolumnwidth-\cvitemwithcommentmainlength}%
  \begin{minipage}[t]{\cvitemwithcommentmainlength}\usebox{\cvitemwithcommentmainbox}\end{minipage}%
  \hfill% fill of \separatorcolumnwidth
  \begin{minipage}[t]{\cvitemwithcommentcommentlength}\raggedleft\small\itshape#4\end{minipage}%
  \par\addvspace{#1}}

%% page layout/margins
\usepackage[top=2.5cm, bottom=2.5cm, left=1.5cm, right=1.5cm]{geometry}

\nopagenumbers{}

%% Babel config for English language
\usepackage[english]{babel}

%% fontspec
\usepackage{fontspec}

\IfFontExistsTF{Linux Libertine}{
  \setmainfont[Ligatures={TeX, Common}, Numbers=Lining, ItalicFont=Linux Libertine]{Linux Libertine}
}{}
\IfFontExistsTF{Linux Libertine O}{
  \setmainfont[Ligatures={TeX, Common}, Numbers=Lining, ItalicFont=Linux Libertine O]{Linux Libertine O}
}{}

%% CTeX
% CJK support, used to show CJK characters in the resume
%
% - fontset=none: disable builtin fontset but instead set the CJK font manually
% - heading=false: disable ctex heading
% - punct=kaiming: use kaiming punctuations styles for CJK
% - scheme=plain: use plain scheme, do not override \normalsize font size
% - space=auto: space settings for CJK characters
%
% ref:
% - http://ctan.mirrorcatalogs.com/language/chinese/ctex/ctex.pdf
\usepackage[UTF8, heading=false, punct=kaiming, scheme=plain, space=auto]{ctex}

\IfFontExistsTF{Noto Serif CJK SC}{
  \setCJKmainfont{Noto Serif CJK SC}
}{}
\IfFontExistsTF{Noto Sans CJK SC}{
  \setCJKsansfont{Noto Sans CJK SC}
}{}

%% line spacing
\usepackage{setspace}
\setstretch{1.125}

%% URL styling - use normal text instead of monospace
\urlstyle{same}

% Auto-underline all links
% Defer until \begin{document} so hyperref is already loaded
\AtBeginDocument{%
  \let\oldhref\href
  \renewcommand{\href}[2]{\oldhref{#1}{\underline{#2}}}%
}

%% Patch \httplink and \httpslink to handle full URLs
% The original moderncv macros always prepend http:// or https:// to URLs.
% This causes issues when the URL already has a protocol (e.g., https://example.com)
% resulting in malformed URLs like https://https://example.com.
% This patch checks if the URL already contains "://" and uses it directly if so.
% Note: We use \str_set:Nx (x-expansion) to expand macros like \@homepage first.
\ExplSyntaxOn
\str_new:N \g_yamlresume_url_str

\RenewDocumentCommand{\httpslink}{O{}m}{
  \str_set:Nx \g_yamlresume_url_str {#2}
  \str_if_in:NnTF \g_yamlresume_url_str {://}
    { \url{#2} }
    { \url{https://#2} }
}

\RenewDocumentCommand{\httplink}{O{}m}{
  \str_set:Nx \g_yamlresume_url_str {#2}
  \str_if_in:NnTF \g_yamlresume_url_str {://}
    { \url{#2} }
    { \url{http://#2} }
}
\ExplSyntaxOff

\name{陳嘉欣}{}
\title{市場營銷經理 | 數碼品牌策略及增長專家}
\phone[mobile]{+852 9123 4567}
\email{ka-yan.chan@example.com}
\homepage{https://www.kayanchan.com}

\address{中國香港，香港島，香港，香港中環德輔道中 88 號}{}{}

\extrainfo{{\small \faLinkedin}\ \href{https://www.linkedin.com/in/kayanchan}{@kayanchan} {} {} {} • {} {} {} 
{\small \faMedium}\ \href{https://medium.com/@kayanchan}{@kayanchan} {} {} {} • {} {} {} 
{\small \faInstagram}\ \href{https://www.instagram.com/kayan.marketing}{@kayan.marketing}}

\begin{document}

\maketitle

\section{簡介}

\cvline{}{擁有超過 8 年市場營銷經驗的專業經理，專注於數碼品牌策略、內容營銷及數據驅動增長。擅長策劃跨渠道營銷活動，提升品牌知名度與客戶轉化率，並善於以數據驗證營銷投資回報。
%
\begin{itemize}
\item 成功領導多個大型營銷項目，平均提升網站流量 120\%，社交媒體互動率 85\%
\item 精通搜尋引擎優化、關鍵字廣告、社交媒體廣告、內容營銷及營銷自動化工具
\item 具備團隊管理經驗，曾帶領 10 人跨職能團隊，協調設計、內容及數據分析工作
\item 熟悉香港及大灣區市場，能制定本地化營銷策略並達成業務目標
\end{itemize}}
\section{教育背景}

\cventry{2013年9月–2017年7月}
        {學士，市場學，成績：3.6}
        {香港中文大學}
        {\url{https://www.cuhk.edu.hk}}
        {}
        {\begin{itemize}
\item 主修市場學，副修心理學，深入理解消費者決策過程
\item 曾參與大學品牌策劃比賽，獲得亞軍
\item 活躍於市場學系學會，協助籌辦業界講座及企業參觀
\end{itemize}
\textbf{課程}：消費者行為、市場研究、品牌管理、數碼營銷、整合營銷傳播、策略管理}
\section{工作經歷}

\cventry{2021年3月至今}
        {高級市場營銷經理}
        {星創科技有限公司}
        {\url{https://www.starcreativetech.hk}}
        {}
        {\begin{itemize}
\item 制定並執行集團整體數碼營銷策略，管理年度營銷預算逾港幣 1,200 萬元
\item 帶領 10 人跨職能團隊，涵蓋內容創作、社交媒體、搜尋引擎優化及數據分析
\item 重新設計品牌網站及內容架構，令自然搜尋流量於 12 個月內增長 120\%
\item 建立營銷自動化流程，將潛在客戶轉化率由 3.1\% 提升至 7.4\%
\item 與產品及銷售團隊緊密協作，主導每季度市場推廣活動的定位與執行
\end{itemize}
\textbf{關鍵字}：品牌策略、數碼營銷、團隊管理、營銷預算}

\cventry{2017年8月–2021年2月}
        {市場營銷主任}
        {匯達廣告有限公司}
        {\url{https://www.huidaglobal.com.hk}}
        {}
        {\begin{itemize}
\item 為零售及金融行業客戶策劃整合營銷活動，平均提升客戶品牌知名度 45\%
\item 負責關鍵字廣告及社交媒體廣告投放，控制每次獲取成本並優化廣告回報率
\item 撰寫中英文營銷文案、新聞稿及社交媒體內容，每月產出超過 30 篇素材
\item 定期製作營銷成效報告，向客戶管理層匯報數據洞察及優化建議
\item 協助培訓新入職同事，建立團隊內部內容審核及品質標準流程
\end{itemize}
\textbf{關鍵字}：整合營銷、廣告投放、內容創作、數據分析}

\cventry{2016年7月–2017年7月}
        {市場推廣助理}
        {藍海市場顧問有限公司}
        {\url{https://www.blueoceanmarketing.hk}}
        {}
        {\begin{itemize}
\item 支援各類線下推廣活動的籌備與執行，包括展覽、路演及產品發佈會
\item 協助管理社交媒體帳號，維護品牌形象並回覆客戶查詢
\item 收集市場及競爭對手資料，整理成月報供團隊參考
\item 與供應商及場地單位協調，確保活動物資及預算符合計劃
\end{itemize}
\textbf{關鍵字}：活動策劃、社交媒體、市場研究}
\section{語言}

\cvline{粵語}{母語或雙語水平 \hfill \textbf{關鍵字}：母語}
\cvline{普通話}{完全專業水平 \hfill \textbf{關鍵字}：普通話水平測試二級甲等}
\cvline{英語}{完全專業水平 \hfill \textbf{關鍵字}：IELTS 7.5}
\section{專業技能}

\cvline{數碼營銷}{專家 \hfill \textbf{關鍵字}：搜尋引擎優化、關鍵字廣告、社交媒體廣告、內容營銷}
\cvline{數據與分析}{高級 \hfill \textbf{關鍵字}：Google Analytics 4、Looker Studio、A/B 測試、數據視覺化}
\cvline{營銷工具}{高級 \hfill \textbf{關鍵字}：HubSpot、Mailchimp、Salesforce、Meta Business Suite}
\cvline{品牌與傳訊}{高級 \hfill \textbf{關鍵字}：品牌定位、公關傳訊、文案撰寫、危機處理}
\cvline{管理與協作}{高級 \hfill \textbf{關鍵字}：團隊領導、預算管理、跨部門協作、供應商談判}
\section{榮譽}

\cventry{2023年11月}
        {香港廣告客戶協會}
        {年度最佳營銷團隊大獎}
        {}
        {}
        {帶領團隊為旗下金融科技產品策劃的整合營銷活動，於年度評選中擊敗逾 60 個參賽項目獲得金獎。}

\cventry{2022年9月}
        {亞太營銷傳播協會}
        {亞太數碼營銷卓越獎（銅獎）}
        {}
        {}
        {憑藉以數據驅動的社交媒體內容策略，成功提升品牌互動率及客戶留存，獲評審團肯定。}
\section{證書}

\cventry{2023年2月}
        {Google}
        {Google Analytics 4 Certification}
        {\url{https://skillshop.exceedlms.com/student/catalog}}
        {}
        {}

\cventry{2021年6月}
        {HubSpot Academy}
        {HubSpot Inbound Marketing Certification}
        {\url{https://academy.hubspot.com/courses}}
        {}
        {}

\cventry{2022年10月}
        {Meta}
        {Meta Certified Digital Marketing Associate}
        {\url{https://www.facebook.com/business/learn/certification}}
        {}
        {}
\section{出版刊物}

\cventry{2023年5月}
        {香港零售品牌的內容營銷突圍策略}
        {亞洲營銷評論}
        {\url{https://www.asiamarketingreview.com}}
        {}
        {\begin{itemize}
\item 分析本地零售品牌如何在內容飽和的環境中建立差異化敘事
\item 提出以客戶旅程為本的內容矩陣框架，並附上實際數據案例
\item 探討人工智能工具在內容生產流程中的應用與限制
\end{itemize}}
\section{項目}

\cventry{2023年1月–2023年8月}
        {為本地環保生活品牌進行的全面品牌重塑及數碼推廣項目}
        {「城巿綠行」品牌重塑計劃}
        {\url{https://www.citygreenwalk.hk}}
        {}
        {\begin{itemize}
\item 重新定義品牌定位與視覺識別，統一線上線下傳訊語言
\item 策劃為期六個月的內容營銷日曆，涵蓋社交媒體、電子報及短片
\item 項目結束時品牌網站流量增長 140\%，社交媒體追蹤人數增加 2.3 倍
\end{itemize}
\textbf{關鍵字}：品牌重塑、內容營銷、可持續發展}

\cventry{2022年3月–2022年11月}
        {針對電商客戶設計的分層自動化再營銷流程與成效追蹤系統}
        {會員再營銷自動化系統}
        {\url{https://www.example.com/crm-automation}}
        {}
        {\begin{itemize}
\item 依據購買行為及互動紀錄將會員劃分為五個生命週期階段
\item 設計對應的電郵及短訊觸發流程，並持續以 A/B 測試優化文案
\item 帶動回購率上升 28\%，並將每次轉化成本降低 19\%
\end{itemize}
\textbf{關鍵字}：營銷自動化、客戶分層、成效追蹤}
\section{興趣愛好}

\cvline{內容創作}{專欄寫作、攝影、短片剪輯}
\cvline{運動與生活}{越野跑、瑜伽、咖啡沖煮}
\section{志願服務}

\cventry{2019年9月至今}
        {義務導師}
        {香港青年營銷協會}
        {\url{https://www.hkyma.org.hk}}
        {}
        {\begin{itemize}
\item 為大專院校學生提供營銷職涯指導，分享行業趨勢及求職準備經驗
\item 協助籌辦年度數碼營銷工作坊，累計服務超過 400 名青年參加者
\item 擔任模擬廣告企劃比賽評審，向參賽隊伍提供專業回饋
\end{itemize}}
    

\cventry{2021年3月–2022年12月}
        {社區推廣義工}
        {綠惜地球}
        {\url{https://www.greenearth.org.hk}}
        {}
        {\begin{itemize}
\item 策劃社區環保推廣活動的社交媒體宣傳，提升公眾參與度
\item 協助撰寫活動新聞稿及宣傳文案，並協調媒體採訪安排
\end{itemize}}
    
\end{document}